\documentclass[12pt,reqno]{amsart}

\usepackage{amsmath,amssymb,amsthm}
\usepackage{booktabs}
\usepackage[colorlinks=true,linkcolor=blue,citecolor=blue,urlcolor=blue]{hyperref}
\usepackage[margin=1.15in]{geometry}
\usepackage{enumitem}
\usepackage{microtype}
\usepackage{longtable}
\usepackage{tcolorbox}

\newtheorem{theorem}{Theorem}[section]
\newtheorem{proposition}[theorem]{Proposition}
\newtheorem{corollary}[theorem]{Corollary}
\theoremstyle{definition}
\newtheorem{definition}[theorem]{Definition}
\theoremstyle{remark}
\newtheorem{remark}[theorem]{Remark}

\newcommand{\ZZ}{\mathbb{Z}}
\newcommand{\QQ}{\mathbb{Q}}
\newcommand{\CC}{\mathbb{C}}
\newcommand{\HH}{\mathbb{H}}
\newcommand{\Gstar}{G^{\!*}}
\newcommand{\varpiup}{\varpi}
\newcommand{\agm}{\operatorname{agm}}
\newcommand{\End}{\operatorname{End}}
\newcommand{\abs}[1]{\lvert #1 \rvert}

\title{Fifty-Two Faces of the Lemniscatic Bridge Constant}

\author{}
\address{}

\subjclass[2020]{33E05, 33B15, 11G05, 11F27}
\keywords{Lemniscate constant, Gamma function, elliptic integrals, theta
  functions, Watson integral, arithmetic-geometric mean, CM elliptic curves,
  Dedekind eta function, Wallis product}

\date{\today}

\begin{document}
\raggedbottom
\allowdisplaybreaks[1]

\begin{abstract}
We compile fifty-two numbered entries concerning the constant
${\Gstar} = \Gamma(\tfrac14)/\Gamma(\tfrac34) \approx 2.9587$,
organised into twelve families spanning the Gamma and Beta functions,
the lemniscate, elliptic integrals, the arithmetic-geometric mean,
Jacobi theta functions, the Dedekind eta function, lattice Green's
functions (Watson integrals), $L$-functions of CM elliptic curves,
hypergeometric series, structural relations among
$\Gstar$, $\varpiup$, and~$\pi$, and modular forms.
The catalogue includes exact representations, identities for related
quantities, algebraic restatements, and one explicitly bounded approximation.
Definitions, normalisations, and the classical identities connecting the
families are stated explicitly.
Among the representations we include the Wallis-type product
$\Gstar = \lim_{N\to\infty} N^{-1/2}\prod_{k=0}^{N-1}(4k+3)/(4k+1)$,
the Dedekind eta evaluation $\Gstar = 2\sqrt{2\pi}\,\eta(i)^{2}$,
and the factorisation of the lemniscate constant as a product
of two residue-class products. These are consequences of classical Gamma,
elliptic, and modular identities; no novelty or priority is claimed.
The collection is a pedagogical dictionary for a normalised lemniscatic
period, not a claim of fifty-two independent derivations.
\end{abstract}

\maketitle


% ================================================================
\section{Definition and notation}\label{sec:def}
% ================================================================

Let $\gamma_{1}=\Gamma(\tfrac14)=3.6256\ldots$ and
$\gamma_{2}=\Gamma(\tfrac12)=\sqrt{\pi}=1.7724\ldots$.

\begin{definition}\label{def:Gstar}
The \emph{lemniscatic bridge constant} is
\begin{equation}\label{eq:def}
  {\Gstar} \;=\; \frac{\gamma_{1}^{2}}{\sqrt{2}\,\gamma_{2}^{2}}
  \;=\; \frac{\sqrt{2}\,\Gamma(\tfrac14)^{2}}{2\pi}
  \;=\; 2.9586751191886388\ldots
\end{equation}
\end{definition}

\begin{tcolorbox}[colback=blue!3!white,colframe=blue!50!black,
  title={The simplest form}]
\begin{equation}\label{eq:simplest}
  \boxed{\;\Gstar \;=\; \frac{\Gamma(\tfrac14)}{\Gamma(\tfrac34)}\;}
\end{equation}
This single ratio of two Gamma values is the identity that makes the
entire collection cohere.  It follows from Euler's reflection formula:
$\Gamma(\tfrac14)\Gamma(\tfrac34) = \pi/\!\sin(\pi/4) = \pi\sqrt{2}$,
so $\Gamma(\tfrac14)/\Gamma(\tfrac34)
= \Gamma(\tfrac14)^{2}/(\pi\sqrt{2})
= \gamma_{1}^{2}/(\sqrt{2}\,\gamma_{2}^{2}) = \Gstar$.
\end{tcolorbox}

\medskip

We write $\varpiup = \gamma_{1}^{2}/(2\sqrt{2}\,\gamma_{2})$
for the lemniscate constant (the half-perimeter of the lemniscate
$r^{2}=\cos 2\theta$) and recall $\pi = \gamma_{2}^{2}$.


% ================================================================
\section{The identities}\label{sec:identities}
% ================================================================

We present fifty-two entries, grouped into twelve families (counting
A$'$ and E$'$ separately). All square roots of positive real quantities
are positive. Except for the approximation E5, the entries are exact
identities, but their left-hand sides are not all $\Gstar$: A10--A12,
A14, C5, F3--F4, I1--I2, and I4 concern related quantities. The summary
table records those targets explicitly.

The entries are not meant to carry equal structural weight.  Some families
give different realisations of the constant: the lemniscate period, elliptic
integral, theta and eta special values, Watson lattice Green's function,
CM-curve $L$-value, and residue-class product.  Other entries are algebraic
corollaries of these realisations, included to make the dictionary complete.
The table in Section~\ref{sec:table} should therefore be read as a catalogue
of representations, not as fifty-two independent derivations.

% ---------------------------------------------------------------
\subsection{Family A: Gamma function and Beta function}\label{sec:A}
% ---------------------------------------------------------------

\begin{align}
  {\Gstar} &= \frac{\gamma_{1}^{2}}{\sqrt{2}\,\gamma_{2}^{2}}
    \tag{A1}\label{A1}\\[5pt]
  {\Gstar} &= \frac{\sqrt{2}\,\Gamma(\tfrac14)^{2}}{2\pi}
    \tag{A2}\label{A2}\\[5pt]
  {\Gstar} &= \frac{B(\tfrac14,\tfrac14)}{\sqrt{2\pi}}
    \tag{A3}\label{A3}\\[5pt]
  {\Gstar} &= \frac{\Gamma(\tfrac14)\,\Gamma(\tfrac14)}
    {\sqrt{2}\,\Gamma(\tfrac12)^{2}}
    \tag{A4}\label{A4}\\[5pt]
  {\Gstar} &= \frac{r}{\sqrt{2}}\,,\qquad
    r = \frac{\gamma_{1}^{2}}{\gamma_{2}^{2}}
    \tag{A5}\label{A5}\\[5pt]
  {\Gstar} &= \frac{\Gamma(\tfrac14)}{\Gamma(\tfrac34)}
    \tag{A6}\label{A6}\\[5pt]
  {\Gstar} &= \frac{\gamma_{2}^{2}\sqrt{2}}{\Gamma(\tfrac34)^{2}}
    \tag{A7}\label{A7}\\[5pt]
  {\Gstar} &= \frac{16\,\Gamma(\tfrac54)^{2}}
    {\sqrt{2}\,\gamma_{2}^{2}}
    \tag{A8}\label{A8}\\[5pt]
  {\Gstar} &= \frac{\bigl(\displaystyle\int_{0}^{\infty}
    t^{-3/4}e^{-t}\,dt\bigr)^{2}}{%
    \sqrt{2}\,\bigl(\displaystyle\int_{0}^{\infty}
    t^{-1/2}e^{-t}\,dt\bigr)^{2}}
    \tag{A9}\label{A9}\\[5pt]
  \sqrt{{\Gstar}} &= \frac{\gamma_{1}}{2^{1/4}\,\gamma_{2}}
    \tag{A10}\label{A10}\\[5pt]
  \frac{1}{{\Gstar}} &= \frac{\Gamma(\tfrac34)}{\Gamma(\tfrac14)}
    \tag{A11}\label{A11}
\end{align}

\noindent
Here $B(a,b)=\Gamma(a)\Gamma(b)/\Gamma(a+b)$ is the Beta function.
\eqref{A6} is the simplest form (see the boxed display above):
a single ratio of two Gamma values.
It follows from the reflection formula
$\Gamma(\tfrac14)\Gamma(\tfrac34)=\pi/\!\sin(\pi/4)=\pi\sqrt{2}
=\sqrt{2}\,\gamma_{2}^{2}$, giving
$\Gamma(\tfrac14)/\Gamma(\tfrac34)
=\Gamma(\tfrac14)^{2}/(\sqrt{2}\,\gamma_{2}^{2})={\Gstar}$.
\eqref{A8} uses $\Gamma(\tfrac54)=\tfrac14\Gamma(\tfrac14)$.
\eqref{A9} substitutes the defining convergent Gamma integrals; it is an
algebraic restatement, not an additional derivation.

% ---------------------------------------------------------------
\subsection{Family \texorpdfstring{A$'$}{A'}: Wallis-type infinite products (the two races)}\label{sec:Aprime}
% ---------------------------------------------------------------

\begin{align}
  \sqrt{\pi} &= \lim_{N\to\infty}\; N^{-1/2}
    \prod_{k=1}^{N}\frac{2k}{2k-1}
    \tag{A12}\label{A12}\\[5pt]
  {\Gstar} &= \lim_{N\to\infty}\; N^{-1/2}
    \prod_{k=0}^{N-1}\frac{4k+3}{4k+1}
    \tag{A13}\label{A13}
\end{align}

\noindent
\eqref{A12} is the classical Wallis product for $\sqrt{\pi}$: the ratio
of even integers over odd integers.  This is \emph{Race~1}---the mod-$2$
partition.

\eqref{A13} is the Wallis-type product for~${\Gstar}$: the ratio
of integers $\equiv 3\pmod{4}$ over integers $\equiv 1\pmod{4}$.
This is \emph{Race~2}---the mod-$4$ partition of all positive odd integers.
For odd \emph{primes}, these residue classes also distinguish inert
and split primes in $\ZZ[i]$; the product is not a product over primes.
Both limits follow from
\[
 \prod_{k=0}^{N-1}\frac{k+a}{k+b}
 =\frac{\Gamma(N+a)\Gamma(b)}{\Gamma(a)\Gamma(N+b)},\qquad
 \frac{\Gamma(N+a)}{\Gamma(N+b)}\sim N^{a-b}.
\]

The lemniscate constant is the product of both races:
\begin{equation}\tag{A14}\label{A14}
  \varpiup \;=\; \frac{{\Gstar}\cdot\sqrt{\pi}}{2}
  \;=\; \frac{1}{2}\,\lim_{N\to\infty}\;N^{-1}\;
  \prod_{k=0}^{N-1}\frac{4k+3}{4k+1}\;\cdot\;
  \prod_{k=1}^{N}\frac{2k}{2k-1}\,.
\end{equation}

\begin{remark}[A grammar of constants]\label{rem:grammar}
This factorisation admits a linguistic metaphor.
$\sqrt{\pi}$ is the \emph{vocabulary}---the even/odd race, the first
partition of the natural numbers.
${\Gstar}$ is the \emph{grammar}---the two odd residue classes modulo four.
$\varpiup$ is the \emph{sentence}---neither race alone,
but both products combined. This is a mnemonic for
$\varpiup=\sqrt\pi\,\Gstar/2$, not a claim of probabilistic independence,
algebraic independence, or a new factorisation theory for constants.
\end{remark}

% ---------------------------------------------------------------
\subsection{Family B: The lemniscate}\label{sec:B}
% ---------------------------------------------------------------

\begin{align}
  {\Gstar} &= \frac{2\varpiup}{\gamma_{2}}
    \tag{B1}\label{B1}\\[5pt]
  {\Gstar} &= \frac{2\varpiup}{\sqrt{\pi}}
    \tag{B2}\label{B2}\\[5pt]
  {\Gstar} &= \frac{\varpiup}{\sqrt{\pi/4}}
    \tag{B3}\label{B3}\\[5pt]
  {\Gstar} &= \frac{4}{\sqrt{\pi}}
    \int_{0}^{1}\frac{dx}{\sqrt{1-x^{4}}}
    \tag{B4}\label{B4}
\end{align}

\noindent
The half-perimeter of the lemniscate $r^{2}=\cos 2\theta$ is~$\varpiup$.
\eqref{B3} writes $\Gstar = \varpiup/\!\sqrt{PF}$ where
$PF = \pi/4$ is the circle-in-square packing fraction.
\eqref{B4} follows by substituting $u=x^4$, which gives the integral
$\tfrac14 B(\tfrac14,\tfrac12)=\varpiup/2$.

% ---------------------------------------------------------------
\subsection{Family C: Elliptic integrals}\label{sec:C}
% ---------------------------------------------------------------

\begin{align}
  {\Gstar} &= \frac{2\sqrt{2}}{\gamma_{2}}\,
    K\!\bigl(\tfrac{1}{\sqrt{2}}\bigr)
    \tag{C1}\label{C1}\\[5pt]
  {\Gstar} &= \frac{2\sqrt{2}}{\sqrt{\pi}}\,
    K\!\bigl(\tfrac{1}{\sqrt{2}}\bigr)
    \tag{C2}\label{C2}\\[5pt]
  {\Gstar} &= \frac{4}{\gamma_{2}}\,K(i)
    \tag{C3}\label{C3}\\[5pt]
  {\Gstar} &= 2\sqrt{\frac{K(1/\!\sqrt{2})}{2\,E(1/\!\sqrt{2})-K(1/\!\sqrt{2})}}
    \tag{C4}\label{C4}\\[5pt]
  {\Gstar}^{2} &= \frac{4\,K(1/\!\sqrt{2})}{2\,E(1/\!\sqrt{2})-K(1/\!\sqrt{2})}
    \tag{C5}\label{C5}
\end{align}

\noindent
Here $K(k)=\int_{0}^{\pi/2}d\theta/\sqrt{1-k^{2}\sin^{2}\theta}$ and
$E(k)=\int_{0}^{\pi/2}\sqrt{1-k^{2}\sin^{2}\theta}\,d\theta$.
The modulus $k=1/\!\sqrt{2}$ is the \emph{lemniscatic modulus}.
It is the unique solution of $j(k)=1728$ in $0<k<1$; complex moduli
related by modular transformations need not be equal.
Here $K$ uses the \emph{modulus} $k$, not the parameter $m=k^2$.
In \eqref{C3}, $K(i)=\int_0^{\pi/2}(1+\sin^2\theta)^{-1/2}d\theta$
is real and positive. The imaginary-modulus transformation gives
$K(i)=K(1/\sqrt2)/\sqrt2$ \cite[19.7.2]{DLMF}.
At the self-complementary modulus, Legendre's relation gives
$K(2E-K)=\pi/2$, proving \eqref{C4}--\eqref{C5} from \eqref{C1}.
The classical evaluation $K(1/\sqrt2)=\gamma_1^2/(4\sqrt\pi)$
connects this family to A \cite[23.5.5]{DLMF}.

% ---------------------------------------------------------------
\subsection{Family D: Arithmetic-geometric mean}\label{sec:D}
% ---------------------------------------------------------------

\begin{align}
  {\Gstar} &= \frac{2\gamma_{2}}{\agm(1,\sqrt{2})}
    \tag{D1}\label{D1}\\[5pt]
  {\Gstar} &= \frac{2\sqrt{\pi}}{\agm(1,\sqrt{2})}
    \tag{D2}\label{D2}\\[5pt]
  {\Gstar} &= \frac{\sqrt{2}\,\gamma_{2}}{\agm(1,1/\!\sqrt{2})}
    \tag{D3}\label{D3}\\[5pt]
  {\Gstar} &= 2\gamma_{2}\prod_{n=0}^{\infty}\frac{2}{1+r_n},
    \qquad r_n=b_n/a_n
    \tag{D4}\label{D4}
\end{align}

\noindent
$\agm(a,b)$ denotes the common limit of the iteration
$a_{n+1}=(a_{n}+b_{n})/2$, $b_{n+1}=\sqrt{a_{n}b_{n}}$.
\eqref{D1} is the classical Gauss identity relating the AGM to
the complete elliptic integral~\cite{BorweinBorwein}.
\eqref{D3} uses the homogeneity $\agm(1,1/\!\sqrt{2})=\agm(1,\sqrt{2})/\sqrt{2}$.
In \eqref{D4}, $(a_0,b_0)=(1,\sqrt2)$ and the preceding AGM iteration
defines every factor. Since $2/(1+r_n)=a_n/a_{n+1}$, its first $N$
factors telescope to $1/a_N$, proving convergence and the identity.
This replaces the undefined cosine product of the earlier version.

% ---------------------------------------------------------------
\subsection{Family E: Jacobi theta functions}\label{sec:E}
% ---------------------------------------------------------------

\begin{align}
  {\Gstar} &= \gamma_{2}\sqrt{2}\;\theta_{3}(e^{-\gamma_{2}^{2}})^{2}
    \tag{E1}\label{E1}\\[5pt]
  {\Gstar} &= \sqrt{2\pi}\;\theta_{3}(e^{-\pi})^{2}
    \tag{E2}\label{E2}\\[5pt]
  {\Gstar} &= \gamma_{2}\sqrt{2}
    \biggl(\sum_{n=-\infty}^{\infty}e^{-\pi n^{2}}\biggr)^{\!2}
    \tag{E3}\label{E3}\\[5pt]
  {\Gstar} &= \gamma_{2}\sqrt{2}\;
    \biggl(1 + 2\sum_{n=1}^{\infty}e^{-\pi n^{2}}\biggr)^{\!2}
    \tag{E4}\label{E4}\\[5pt]
  {\Gstar} &\approx \gamma_{2}\sqrt{2}\;
    (1 + 2\,e^{-\pi})^{2}
    \quad\text{(approximation; error bounded below)}
    \tag{E5}\label{E5}
\end{align}

\noindent
$\theta_{3}(q)=\sum_{n\in\ZZ}q^{n^{2}}$ is the Jacobi theta function.
For positive $t$, Poisson summation gives
$\theta_{3}(e^{-\pi t})=t^{-1/2}\theta_{3}(e^{-\pi/t})$.
The map $t\mapsto1/t$ fixes precisely $t=1$, hence $q_0=e^{-\pi}$.
The relation $K(k)=\pi\theta_3(q)^2/2$ with
$q=\exp(-\pi K(\sqrt{1-k^2})/K(k))$ \cite[20.9.2]{DLMF}
and the evaluation in Family C give
$\theta_{3}(e^{-\pi})=\gamma_{1}/(\sqrt2\,\gamma_{2}^{3/2})$.
This is also a special case of the CM-period evaluations organised by
the Chowla--Selberg formula~\cite{ChowlaSelberg}.

For a rigorous bound on E5, put $q=e^{-\pi}$, $A=1+2q$ and
$B=2q^4/(1-q^5)$. Since successive exponents in
$\sum_{n=2}^\infty q^{n^2}$ differ by at least five,
$0<\theta_3(q)-A\leq B$. Consequently
\[
 0<\Gstar-\sqrt{2\pi}\,A^2
 \leq\sqrt{2\pi}(2AB+B^2)<3.7989\times10^{-5}.
\]
The approximation is $2.9586371311\ldots$, not an exact identity.

% ---------------------------------------------------------------
\subsection{Family \texorpdfstring{E$'$}{E'}: Dedekind eta function}\label{sec:Eprime}
% ---------------------------------------------------------------

\begin{align}
  {\Gstar} &= 2\gamma_{2}\sqrt{2}\;\eta(i)^{2}
    \tag{E6}\label{E6}\\[5pt]
  {\Gstar} &= 2\sqrt{2\pi}\;\eta(i)^{2}
    \tag{E7}\label{E7}
\end{align}

\noindent
$\eta(\tau)=e^{\pi i\tau/12}\prod_{n=1}^{\infty}(1-e^{2\pi in\tau})$
is the Dedekind eta function.
At $\tau=i$: $\eta(i)=\theta_3(e^{-\pi})/\sqrt2
=\gamma_{1}/(2\gamma_{2}^{3/2})$.
For example, this follows from
$2\eta(\tau)^3=\theta_2(\tau)\theta_3(\tau)\theta_4(\tau)$,
$\theta_2(i)=\theta_4(i)$ and
$\theta_3(i)^4=\theta_2(i)^4+\theta_4(i)^4$, using positive values.
Then $2\gamma_{2}\sqrt{2}\cdot\eta(i)^{2}
=2\gamma_{2}\sqrt{2}\cdot\gamma_{1}^{2}/(4\gamma_{2}^{3})
=\gamma_{1}^{2}/(\sqrt{2}\,\gamma_{2}^{2})={\Gstar}$.

% ---------------------------------------------------------------
\subsection{Family F: Lattice sums (Watson integrals)}\label{sec:F}
% ---------------------------------------------------------------

\begin{align}
  {\Gstar} &= \sqrt{2\gamma_{2}^{2}\,W_{3}}
    \tag{F1}\label{F1}\\[5pt]
  {\Gstar} &= \sqrt{2\pi\,W_{3}}
    \tag{F2}\label{F2}\\[5pt]
  {\Gstar}^{2} &= 2\gamma_{2}^{2}\,W_{3}
    \tag{F3}\label{F3}\\[5pt]
  {\Gstar}^{2} &= \frac{1}{4\pi^{2}}
    \iiint_{[-\pi,\pi]^{3}}
    \frac{dk_{1}\,dk_{2}\,dk_{3}}{%
    1-\cos k_{1}\cos k_{2}\cos k_{3}}
    \tag{F4}\label{F4}
\end{align}

\noindent
$W_{3}=\gamma_{1}^{4}/(4\gamma_{2}^{6})$ is the Watson integral for
the body-centred cubic (BCC) lattice~\cite{Watson}.
Our convention is
\[
 W_3=\frac{1}{(2\pi)^3}\iiint_{[-\pi,\pi]^3}
 \frac{d^3k}{1-\cos k_1\cos k_2\cos k_3}
 =\frac{4K(1/\sqrt2)^2}{\pi^2}.
\]
This is the normalised Fourier-cell integral for the BCC step symbol
$\cos k_1\cos k_2\cos k_3$. Its isolated singularities are locally
integrable in three dimensions. The integral is the increasing limit
with the denominator replaced by $1-z\cos k_1\cos k_2\cos k_3$
after averaging paired signs, as $z\uparrow1$.
Expanding at $0\leq z<1$ and averaging the even powers gives
$\sum_{n\geq0}\binom{2n}{n}^3z^{2n}/64^n$; Watson's evaluation at
$z=1$ gives the stated value. The FCC symbol
$\cos k_1\cos k_2+\cos k_2\cos k_3+\cos k_3\cos k_1$
belongs to a different integral and cannot be substituted in F4.

% ---------------------------------------------------------------
\subsection{Family G: A fixed CM elliptic curve and its \texorpdfstring{$L$}{L}-value}\label{sec:G}
% ---------------------------------------------------------------

\begin{align}
  {\Gstar} &= \frac{8\,L(E,1)}{\gamma_{2}}
    \tag{G1}\label{G1}\\[5pt]
  {\Gstar} &= \frac{8\,L(E,1)}{\sqrt{\pi}}
    \tag{G2}\label{G2}\\[5pt]
  {\Gstar} &= \frac{8}{\sqrt\pi}\;L(E,1)
    \tag{G3}\label{G3}\\[5pt]
  {\Gstar} &= \frac{2\,\Omega_{+}}{\gamma_{2}}
    \tag{G4}\label{G4}
\end{align}

\noindent
$E\colon y^{2}=x^{3}-x$ is the CM elliptic curve with $j$-invariant~$1728$
and $\End_{\overline{\QQ}}(E)\cong\ZZ[i]$ (the full CM endomorphisms
are not all defined over $\QQ$). Here $L(E,s)$ is the uncompleted
Hasse--Weil $L$-function, with central point $s=1$ and $a_1=1$;
the conductor is $32$ and the curve is LMFDB $32.\mathrm{a}3$
(Cremona $32\mathrm{a}2$) \cite{LMFDB}.
For the N\'eron differential $\omega=dx/(2y)$, define
\[
 \Omega_+=\int_1^\infty\frac{dx}{\sqrt{x^3-x}}=\varpiup,
 \qquad \int_{E(\mathbb R)}|\omega|=2\varpiup.
\]
The first equality with $\varpiup$ follows by $x=t^{-2}$.
Li, Long and Tu \cite[Lemma 3.1]{LiLongTu} prove directly that
\[
 L(E,1)=\tfrac18 B(\tfrac12,\tfrac14)
       =\Omega_+/4=\varpiup/4.
\]
G1--G3 are algebraic restatements of this fixed-curve evaluation.
It is consistent with the rank-zero BSD quotient
$(2\varpiup)\cdot1\cdot2/4^2=\varpiup/4$, using both real
components, Tamagawa number $c_2=2$, torsion order four and trivial
Tate--Shafarevich group. No general BSD theorem or claim about twists
is needed for these entries. The exact value can instead be obtained
by integrating the weight-two form
$f(\tau)=\eta(4\tau)^2\eta(8\tau)^2$ associated to this curve.
In particular its Mellin integral is
$L(E,1)=2\pi\int_0^\infty f(it)\,dt$; the numerical verification
also evaluates its coefficients independently of the Gamma formula.

% ---------------------------------------------------------------
\subsection{Family H: Hypergeometric and series representations}\label{sec:H}
% ---------------------------------------------------------------

\begin{align}
  {\Gstar} &= \sqrt{2\pi}\;
    {}_{2}F_{1}\!\bigl(\tfrac12,\tfrac12;\,1;\,\tfrac12\bigr)
    \tag{H1}\label{H1}\\[5pt]
  {\Gstar} &= \gamma_{2}\sqrt{2}\;
    \biggl(\sum_{n=0}^{\infty} r_{2}(n)\,e^{-\gamma_{2}^{2}n}\biggr)
    \tag{H2}\label{H2}\\[5pt]
  {\Gstar} &= \gamma_{2}\sqrt{2}\;
    \biggl(1 + 4\sum_{\substack{d\geq 1\\d\;\mathrm{odd}}}
    \frac{\chi_{-4}(d)\,e^{-\gamma_{2}^{2}d}}
    {1-e^{-\gamma_{2}^{2}d}}\biggr)
    \tag{H3}\label{H3}\\[5pt]
  {\Gstar} &= \gamma_{2}\sqrt{2}\;
    \prod_{n=1}^{\infty}
    (1-e^{-2\gamma_{2}^{2}n})^{2}\,
    (1+e^{-\gamma_{2}^{2}(2n-1)})^{4}
    \tag{H4}\label{H4}\\[5pt]
  {\Gstar} &= \frac{16}{\sqrt{2}\,\gamma_{2}^{2}}
    \biggl(\int_{0}^{\infty}e^{-t^{4}}\,dt\biggr)^{\!2}
    \tag{H5}\label{H5}
\end{align}

\noindent
\eqref{H1} uses $K(1/\!\sqrt{2})=(\pi/2)\;
{}_{2}F_{1}(\tfrac12,\tfrac12;1;\tfrac12)$.
\eqref{H2} expresses $\theta_{3}^{2}$ via the arithmetic function
$r_{2}(n)=\#\{(a,b)\in\ZZ^{2}:a^{2}+b^{2}=n\}$, with $r_2(0)=1$.
\eqref{H3} is the Lambert-series form using the non-principal character
$\chi_{-4}(d)=\bigl(\frac{-4}{d}\bigr)$: Jacobi's divisor formula
$r_2(n)=4\sum_{d\mid n}\chi_{-4}(d)$ for $n\geq1$ and an absolutely
convergent rearrangement give H3 from H2.
\eqref{H4} is the Jacobi triple-product formula applied to $\theta_{3}^{2}$.
\eqref{H5} uses $\int_{0}^{\infty}e^{-t^{4}}\,dt=\Gamma(\tfrac54)=\tfrac14\Gamma(\tfrac14)$.

% ---------------------------------------------------------------
\subsection{Family I: Structural relations among \texorpdfstring{$\Gstar$, $\varpiup$, and $\pi$}{G*, varpi, and pi}}\label{sec:I}
% ---------------------------------------------------------------

\begin{align}
  \pi &= \frac{4\,\varpiup^{2}}{{\Gstar}^{2}}
    \tag{I1}\label{I1}\\[5pt]
  \varpiup &= \frac{{\Gstar}\,\gamma_{2}}{2}
    \tag{I2}\label{I2}\\[5pt]
  {\Gstar} &= \frac{4\,\varpiup^{2}}{{\Gstar}\,\gamma_{2}^{2}}
    \tag{I3}\label{I3}\\[5pt]
  {\Gstar}^{3} &= \frac{16\,\varpiup^{4}}
    {{\Gstar}\,\gamma_{2}^{4}}
    \tag{I4}\label{I4}
\end{align}

\noindent
\eqref{I1} is the \emph{triad identity}: the circular constant expressed
as a ratio of lemniscatic objects.
\eqref{I3} is the self-referential form of~\eqref{I1}.
\eqref{I4} is another algebraic rearrangement and still contains
$\Gstar$ on its right-hand side. Neither I3 nor I4 is an independent
evaluation method or a new relation among independent constants.

% ---------------------------------------------------------------
\subsection{Family J: Modular forms}\label{sec:J}
% ---------------------------------------------------------------

\begin{align}
  {\Gstar} &= \sqrt{2\pi}\;\frac{\theta_{3}(e^{-\pi})^{4}}
    {\theta_{3}(e^{-\pi})^{2}}
    = \sqrt{2\pi}\;\theta_{3}(e^{-\pi})^{2}
    \tag{J1}\label{J1}
\end{align}

\noindent
\eqref{J1} is deliberately a restatement of~\eqref{E2}.
Write $\Theta(\tau)=\sum_{n\in\ZZ}e^{\pi i n^2\tau}$, so the
theta nome is $q=e^{\pi i\tau}$, not the eta-product nome
$e^{2\pi i\tau}$. Its square has weight one with a multiplier on
the theta subgroup generated by $S:\tau\mapsto-1/\tau$ and
$T^2:\tau\mapsto\tau+2$:
\[
 \Theta(-1/\tau)^2=(-i\tau)\Theta(\tau)^2,\qquad
 \Theta(\tau+2)^2=\Theta(\tau)^2.
\]
These are the theta transformation laws \cite[20.7]{DLMF}.
Then $\Gstar=\sqrt{2\pi}\,\Theta(i)^2$.
The unique fixed point of the specific involution $S$ in $\HH$ is $i$;
this does not assert uniqueness of elliptic points on the modular
quotient (which also has the order-three point $e^{2\pi i/3}$).
The nonholomorphic factor $\sqrt{\operatorname{Im}\tau}$ would not
define a holomorphic weight-one modular form. $\Gstar$ is a scalar
special value associated with the involution, not itself an operator.


% ================================================================
\clearpage
\section{Summary table}\label{sec:table}
% ================================================================

The following table lists all fifty-two entries. Unless an alternative
left-hand side is displayed, the expression equals $\Gstar$; E5 is an
approximation only. In rows C4--C5, $K$ and $E$ are evaluated at
$k=1/\sqrt2$. The
\textbf{Source} column gives either the original reference or the
identity from which the entry follows by algebraic manipulation.

\begin{center}
\begin{longtable}{clll}
\toprule
\textbf{ID} & \textbf{Expression (target if different)} & \textbf{Family} & \textbf{Source} \\
\midrule
\endfirsthead
\toprule
\textbf{ID} & \textbf{Expression (target if different)} & \textbf{Family} & \textbf{Source} \\
\midrule
\endhead
A1 & $\gamma_{1}^{2}/(\sqrt{2}\,\gamma_{2}^{2})$ & Gamma & Definition \\
A2 & $\sqrt{2}\,\Gamma(\frac14)^{2}/(2\pi)$ & Gamma & Classical \\
A3 & $B(\frac14,\frac14)/\sqrt{2\pi}$ & Beta & Classical \\
A4 & $\Gamma(\frac14)^{2}/(\sqrt{2}\,\Gamma(\frac12)^{2})$ & Gamma & \eqref{A1} \\
A5 & $r/\sqrt{2}$, \; $r=\gamma_1^2/\gamma_2^2$ & Gamma ratio & Algebra \\
A6 & $\Gamma(\frac14)/\Gamma(\frac34)$ & Gamma ratio & Reflection \\
A7 & $\gamma_2^2\sqrt{2}/\Gamma(\frac34)^2$ & Gamma & Reflection \\
A8 & $16\,\Gamma(\frac54)^2/(\sqrt{2}\,\gamma_2^2)$ & Gamma & Recurrence \\
A9 & Ratio of Laplace integrals & Pure integral & Definition \\
A10 & $\sqrt{{\Gstar}}=\gamma_1/(2^{1/4}\gamma_2)$ & Square root & Algebra \\
A11 & $1/{\Gstar}=\Gamma(\frac34)/\Gamma(\frac14)$ & Reciprocal & \eqref{A6} \\
A12 & $\sqrt\pi = \lim N^{-1/2}\prod(2k)/(2k{-}1)$ & Race 1 & Wallis~\cite{Wallis} \\
A13 & ${\Gstar} = \lim N^{-1/2}\prod(4k{+}3)/(4k{+}1)$ & Race 2 & Gamma ratio \\
A14 & $\varpiup = \tfrac12\cdot\text{Race\,1}\times\text{Race\,2}$ & Both races & \eqref{A12}$+$\eqref{A13} \\
\midrule
B1 & $2\varpiup/\gamma_{2}$ & Lemniscate & Definition \\
B2 & $2\varpiup/\sqrt{\pi}$ & Lemniscate & \eqref{B1} \\
B3 & $\varpiup/\sqrt{\pi/4}$ & Lemniscate & \eqref{B2} \\
B4 & $(4/\sqrt{\pi})\int_{0}^{1}(1-x^{4})^{-1/2}\,dx$ & Integral & Beta integral \\
\midrule
C1 & $2\sqrt{2}\,K(1/\!\sqrt{2})/\gamma_{2}$ & Elliptic $K$ & Legendre \\
C2 & $2\sqrt2 K(1/\!\sqrt{2})/\sqrt{\pi}$ & Elliptic $K$ & \eqref{C1} \\
C3 & $4K(i)/\gamma_{2}$ & Imaginary modulus & \cite[19.7.2]{DLMF} \\
C4 & $2\sqrt{K/(2E-K)}$ & Elliptic $E,K$ & Legendre \\
C5 & ${\Gstar}^2=4K/(2E-K)$ & Elliptic squared & \eqref{C4} \\
\midrule
D1 & $2\gamma_{2}/\agm(1,\sqrt{2})$ & AGM & \cite[19.8.5]{DLMF} \\
D2 & $2\sqrt{\pi}/\agm(1,\sqrt{2})$ & AGM & \eqref{D1} \\
D3 & $\sqrt2\,\gamma_{2}/\agm(1,1/\!\sqrt{2})$ & AGM & \eqref{D1} \\
D4 & $2\gamma_2\prod_{n\geq0}2/(1+r_n)$ & AGM product & Telescoping \\
\midrule
E1 & $\gamma_{2}\sqrt{2}\,\theta_{3}(e^{-\pi})^{2}$ & Theta & Chowla--Selberg \\
E2 & $\sqrt{2\pi}\,\theta_{3}(e^{-\pi})^{2}$ & Theta & \eqref{E1} \\
E3 & $\gamma_2\sqrt{2}\,(\sum e^{-\pi n^2})^2$ & Theta series & \eqref{E1} \\
E4 & $\gamma_2\sqrt{2}\,(1+2\sum e^{-\pi n^2})^2$ & Theta half-sum & \eqref{E3} \\
E5 & $\approx\gamma_2\sqrt{2}\,(1+2e^{-\pi})^2$ & Approximation & Bounded tail \\
E6 & $2\gamma_2\sqrt{2}\,\eta(i)^2$ & Dedekind $\eta$ & Chowla--Selberg \\
E7 & $2\sqrt{2\pi}\,\eta(i)^2$ & Dedekind $\eta$ & \eqref{E6} \\
\midrule
F1 & $\sqrt{2\gamma_{2}^{2}W_{3}}$ & Watson & Watson~\cite{Watson} \\
F2 & $\sqrt{2\pi W_{3}}$ & Watson & \eqref{F1} \\
F3 & ${\Gstar}^{2}=2\gamma_{2}^{2}W_{3}$ & Watson squared & \eqref{F1} \\
F4 & ${\Gstar}^2=$ BCC Fourier-cell integral & Watson explicit & Watson~\cite{Watson} \\
\midrule
G1 & $8L(E,1)/\gamma_{2}$ & $L$-function & \cite[Lemma 3.1]{LiLongTu} \\
G2 & $8L(E,1)/\sqrt{\pi}$ & $L$-function & \eqref{G1} \\
G3 & $(8/\sqrt\pi)\,L(E,1)$ & $L$-function & \eqref{G1} \\
G4 & $2\Omega_{+}/\gamma_{2}$ & Period & Integral substitution \\
\midrule
H1 & $\sqrt{2\pi}\,{}_{2}F_{1}(\frac12,\frac12;1;\frac12)$ & Hypergeometric & \cite[19.5.1]{DLMF} \\
H2 & $\gamma_2\sqrt{2}\sum r_{2}(n)\,e^{-\pi n}$ & Lattice count & Jacobi \\
H3 & Lambert series with $\chi_{-4}$ & Dirichlet char.\ & Classical \\
H4 & Jacobi triple product for $\theta_3^2$ & Product & Jacobi 1829 \\
H5 & $\frac{16}{\sqrt{2}\pi}(\int_0^\infty e^{-t^4}dt)^2$ & Quartic integral & \eqref{A8} \\
\midrule
I1 & $\pi=4\varpiup^{2}/{\Gstar}^{2}$ (triad) & Triad & \eqref{B1} \\
I2 & $\varpiup={\Gstar}\gamma_{2}/2$ & Triad & \eqref{B1} \\
I3 & ${\Gstar}=4\varpiup^2/({\Gstar}\gamma_2^2)$ & Triad & \eqref{I1} \\
I4 & ${\Gstar}^{3}=16\varpiup^4/({\Gstar}\gamma_2^4)$ & Triad & \eqref{I1}$+$\eqref{I3} \\
\midrule
J1 & $\sqrt{2\pi}\,\theta_3(e^{-\pi})^2$ (modular) & Modular form & Chowla--Selberg \\
\bottomrule
\end{longtable}
\end{center}


% ================================================================
\section{Convergence of the Wallis products}\label{sec:convergence}
% ================================================================

The limits in~\eqref{A12} and~\eqref{A13} are exact identities.
Their finite truncations are approximations. Using $N$ factors in
each product makes their different convergence rates transparent.

\begin{theorem}[Convergence rate]\label{thm:rate}
Let $S_N^{(\pi)}=N^{-1/2}\prod_{k=1}^N(2k)/(2k-1)$ and
$S_N^{(G)}=N^{-1/2}\prod_{k=0}^{N-1}(4k+3)/(4k+1)$. Then
\[
 \log\frac{S_N^{(\pi)}}{\sqrt\pi}
   =\frac1{8N}-\frac1{192N^3}+O(N^{-5}),\qquad
 \log\frac{S_N^{(G)}}{\Gstar}
   =\frac1{64N^2}-\frac5{2048N^4}+O(N^{-6}).
\]
In particular the raw errors are $O(N^{-1})$ and $O(N^{-2})$,
respectively. The corrected approximations
$S_N^{(\pi)}e^{-1/(8N)}$ and $S_N^{(G)}e^{-1/(64N^2)}$
have errors $O(N^{-3})$ and $O(N^{-4})$, respectively.
\end{theorem}

\begin{proof}
The finite Gamma-product identity in Family A$'$ gives exactly
\[
 \frac{S_N^{(\pi)}}{\sqrt\pi}
 =\frac{\Gamma(N+1)}{\sqrt N\,\Gamma(N+1/2)},\qquad
 \frac{S_N^{(G)}}{\Gstar}
 =\frac{\Gamma(N+3/4)}{\sqrt N\,\Gamma(N+1/4)}.
\]
For fixed $a,b$, subtract the shifted log-Gamma expansions
\cite[5.11.8]{DLMF}:
\[
 \log\frac{\Gamma(N+a)}{\Gamma(N+b)}
 \sim(a-b)\log N+
 \sum_{r\geq1}\frac{(-1)^{r+1}
 [B_{r+1}(a)-B_{r+1}(b)]}{r(r+1)N^r},
\]
where $B_r(x)$ are the Bernoulli polynomials. Substitution of
$(a,b)=(1,1/2)$ and $(3/4,1/4)$ gives the displayed coefficients.
Subtracting the leading logarithmic correction and exponentiating
proves the error orders.
\end{proof}

\begin{remark}[Rates and computation]\label{rem:digits}
An error of order $N^{-p}$ gains asymptotically $p$ decimal digits per
factor of ten in $N$, provided arithmetic precision is sufficient.
Thus the two raw sequences have different rates, and so do the
specified corrected sequences. Additional Stirling terms can further
accelerate them; any numerical guarantee requires a truncation bound
and an arithmetic precision policy. Replacing the prefactor by the
\emph{exact} reciprocal Gamma ratio gives an exact finite identity,
not an approximation with an improved error order. No identical-cost
or large-workstation-feasibility claim follows from these formulas.
\end{remark}

\begin{remark}[Exactness]\label{rem:exact}
The Wallis product for $\sqrt{\pi}$ is not an approximation to $\sqrt{\pi}$;
it \emph{equals} $\sqrt{\pi}$ as a limit.
Likewise, the product~\eqref{A13} for ${\Gstar}$
\emph{equals}~${\Gstar}$.  Both are exact definitions of their
respective constants as limits of finite products over residue classes.
The first race (evens versus odds) defines~$\sqrt\pi$.
The second race ($1\bmod 4$ versus $3\bmod 4$) defines~${\Gstar}$.
\end{remark}

The identity ${\Gstar}^{2}=8K(1/\sqrt2)^2/\pi$ is an exact consequence
of C1. Numerical agreement is a useful implementation check, but does
not supply an independent proof or certify an unevaluated product limit.


% ================================================================
\section{Numerical verification}\label{sec:numerical}
% ================================================================

The defining Gamma ratio evaluates to
\[
  {\Gstar} = 2.9586751191886388\ldots
\]
using arbitrary-precision arithmetic. The companion catalogue and
verification script distinguish exact identities from finite
approximations and compare each entry to its \emph{stated target}.
For limits and integrals, the method and tolerance matter: checking a
closed form is not equivalent to evaluating an improper triple
integral to the same precision. E5 has the explicit error bound above.
The reproducible entry checks are in
\texttt{dissemination/interactive/gstar-introduction/}\allowbreak
\texttt{verify\_catalogue.py}, with a per-entry report in
\texttt{CATALOGUE\_VERIFICATION.json} in the same directory.

The earlier blanket claims that all fifty-two entries equal $\Gstar$
to double precision, agree to $200$ decimals, and are machine-verified
in Lean are withdrawn. In particular, the relevant declarations in
\texttt{lean/FTD/Axioms.lean} use proposition \texttt{True}; they record
citations and do not formalise these special-function identities.
The proofs and primary references in this paper provide the mathematical
basis; computation supplies reproducible checks, not a formal proof.
For broader background on related constants, see Finch~\cite{Finch}.

\medskip

The value of this catalogue is the network of classical connections,
with their normalisations visible. The entry count includes repeated
forms and identities for related constants; it is not a count of
independent discoveries.


% ================================================================
% REFERENCES
% ================================================================

\begin{thebibliography}{10}

\bibitem{BorweinBorwein}
J.~M.~Borwein and P.~B.~Borwein,
\emph{Pi and the AGM: A Study in Analytic Number Theory and Computational
Complexity}, Wiley, New York, 1987.

\bibitem{ChowlaSelberg}
S.~Chowla and A.~Selberg,
On Epstein's zeta-function,
\emph{J.~Reine Angew.\ Math.}\ \textbf{227} (1967), 86--110.

\bibitem{DLMF}
NIST Digital Library of Mathematical Functions,
\url{https://dlmf.nist.gov/}. Specific formulas used here:
5.11.8 (shifted log-Gamma expansion), 19.5.1 (hypergeometric expansion),
19.7.1--2 (Legendre relation and imaginary modulus), 19.8.5 (AGM),
20.7 (theta transformations), 20.9.2 (theta and elliptic integrals),
and 23.5.5 (lemniscatic special value). Accessed 24 September 2026.

\bibitem{Finch}
S.~R.~Finch,
\emph{Mathematical Constants},
Encyclopedia of Mathematics and its Applications~\textbf{94},
Cambridge University Press, 2003.

\bibitem{LiLongTu}
W.-C.~W.~Li, L.~Long and F.-T.~Tu,
Computing special $L$-values of certain modular forms with complex
multiplication, \emph{SIGMA} \textbf{14} (2018), 090, 32 pp.,
Lemma 3.1. \url{https://doi.org/10.3842/SIGMA.2018.090}.

\bibitem{LMFDB}
The LMFDB Collaboration, The $L$-functions and Modular Forms Database,
elliptic curve $32.\mathrm{a}3$ and modular form $32.2.\mathrm{a.a}$,
\url{https://www.lmfdb.org/EllipticCurve/Q/32/a/3} and
\url{https://www.lmfdb.org/ModularForm/GL2/Q/holomorphic/32/2/a/a/}.
Accessed 24 September 2026.

\bibitem{Wallis}
J.~Wallis,
\emph{Arithmetica Infinitorum}, Oxford, 1656.

\bibitem{Watson}
G.~N.~Watson,
Three triple integrals,
\emph{Quart.\ J.\ Math.\ Oxford}\ \textbf{10} (1939), 266--276.

\end{thebibliography}

\end{document}
